\subsubsection{Core Functional Analytic Theorems}

The basic theorems of functional analysis explain why complete normed spaces behave rigidly under linear maps. Baire category supplies the hidden engine: pointwise control, surjectivity, and closed graphs become uniform operator statements.

\begin{theorem} \label{thm:hahn-banach-extension} \hlt{Hahn--Banach Extension Theorem}\\
Let $\mathcal{X}$ be a real normed space, let $M\subseteq\mathcal{X}$ be a linear subspace, and let $p:\mathcal{X}\to\R$ be sublinear. If $\varphi:M\to\R$ is linear and $\varphi(x)\leq p(x)$ for all $x\in M$, then $\varphi$ extends to a linear functional $\Phi:\mathcal{X}\to\R$ with $\Phi(x)\leq p(x)$ for all $x\in\mathcal{X}$. In the normed complex form, every bounded functional on a subspace extends with the same norm.
\end{theorem}

\begin{theorem} \label{thm:uniform-boundedness-principle} \hlt{Uniform Boundedness Principle}\\
Let $\mathcal{X}$ be Banach and let $\mathcal{Y}$ be normed. If $\mathcal{F}\subseteq\LL(\mathcal{X},\mathcal{Y})$ is pointwise bounded, meaning $\sup_{T\in\mathcal{F}}\norm{Tx}<\infty$ for every $x\in\mathcal{X}$, then $\sup_{T\in\mathcal{F}}\normop{T}<\infty$.
\end{theorem}

\begin{theorem} \label{thm:open-mapping-theorem} \hlt{Open Mapping Theorem}\\
A surjective bounded linear map $T:\mathcal{X}\to\mathcal{Y}$ between Banach spaces is open. Equivalently, $T$ maps neighborhoods of $0$ onto neighborhoods of $0$.
\end{theorem}

\begin{theorem} \label{thm:closed-graph-theorem} \hlt{Closed Graph Theorem}\\
Let $\mathcal{X}$ and $\mathcal{Y}$ be Banach spaces. A linear map $T:\mathcal{X}\to\mathcal{Y}$ is bounded if and only if its graph
\[
\Gamma(T)\defeq\{(x,Tx):x\in\mathcal{X}\}
\]
is closed in $\mathcal{X}\times\mathcal{Y}$.
\end{theorem}
